

\section{Basics}

\begin{frame}{Qu'est-ce que \mongodb/,?}

%\frametitle{Qu'est-ce que \mongodb/,?}

\mongodb/, un système "NoSQL", l'un des plus populaires

\begin{redbullet}
  \item fait partie des NoSQL dits ``documentaires''

  \item s'appuie sur un modèle de données semi-structuré (encodage JSON)\,;
  \item pas de schéma (complète flexibilité)\,;
  \item un \alert{langage d'interrogation} original (et spécifique)\,;
   \item pas (ou très peu) de \alert{support transactionnel}.
   \end{redbullet}
   \vfill
   
Construit dès l'origine comme un système \alert{scalable} et \alert{distribué}.

\begin{redbullet}
 \item distribution par partionnement (\textit{sharding})\,;
  \item technique adoptée\,: découpage par intervalles (type BigTable, Google)\,;
  \item tolérance aux pannes par réplication.
\end{redbullet}   

\end{frame}


\begin{frame}
\frametitle{Objectifs de ce cours}


\alert{Découvrir un système NoSQL} $\Rightarrow$ session interactive commentée.

\vfill

\alert{Modélisation}. Comment modélise-t-on une application\,? 
Avantages/inconvénients vs. un système relationnel.

\vfill

Un aperçu du \alert{langage d'interrogation} de \mongodb/.

\end{frame}


\section{Session découverte\,!}


\begin{frame}
\frametitle{Le minimum à savoir}

Vous connaissez déjà JSON\,?

\begin{redbullet}
 \item Système client/serveur\,; le serveur est \texttt{mongod}
 \item Open-source, multi plate-formes\,: installez-le chez vous\,!
  \item Un client en mode terminal\,: \texttt{mongo}\,; langage Javascript\,;
  \item D'autres clients plus ergonomiques, dont \texttt{RockMongo}, une interface Web.
\end{redbullet}

Conçu et développé par 10gen.

\end{frame}

\begin{frame}
\frametitle{Découvrons (avec le client \texttt{mongo})}

On se place dans une \alert{base} \texttt{nfe204} (cd. MySQL)

\smallskip

\texttt{use nfe024}

\smallskip

On insère un \alert{document} JSON dans une \alert{collection} 

\smallskip

\texttt{> db.test.insert (\{"nom": "nfe024"\})}

\smallskip

On affiche le contenu de la collection

\smallskip

\texttt{> db.test.find() }

On insère n'importe quel autre document (pas de \alert{schéma})

\smallskip

\texttt{> db.test.insert (\{"produit": "Grulband", prix: 230, enStock: true\})}

\smallskip

On peut donner un identifiant explicite\,:

\smallskip

\texttt{> db.test.insert (\{\_id: "1", "produit": "Kramölk", prix: 10, enStock: true\})}

\smallskip

\texttt{> db.test.count()}

\end{frame}



\begin{frame}[fragile]
\frametitle{Quelques documents plus sérieux}

Exemple

\smallskip


\begin{minted}[baselinestretch=.9,
               fontsize=\small,
               framesep=2mm]{javascript}
{
  "_id": "movie:100",
  "title": "The Social network", 
  "summary": "On a fall night in 2003, Harvard undergrad and 
         programming genius Mark Zuckerberg sits down at his  
         computer and heatedly begins working on a new idea. (...)",
  "year": 2010, 
  "director": {"last_name": "Fincher",
                    "first_name": "David"},
  "actors": [ 
        {"first_name": "Jesse", "last_name": "Eisenberg"}, 
        {"first_name": "Rooney", "last_name": "Mara"}
    ]
}
\end{minted}

\end{frame}


\begin{frame}[fragile]
\frametitle{Aperçu du langage de requêtes}

L'équivalent du \texttt{select-from-where}, \alert{sur une
seule collection} (pas de \alert{jointure}).

\vfill


\begin{smallverbatim}
db.movies.find ({"_id": "movie:2"})
\end{smallverbatim}

\vfill

Avec des \alert{expressions de chemin} (simples!)

\vfill


\begin{smallverbatim}
db.test.find ({"director.last_name": "Scott"} )
\end{smallverbatim}

\vfill
Fonctionne \alert{aussi} pour les acteurs (même si c'est un \alert{tableau}).

\vfill
\begin{smallverbatim}
db.test.find ({"actors.last_name": "Weaver"})
\end{smallverbatim}

\vfill
Mise à jour\,:

\vfill

\begin{smallverbatim}
db.test.update ({"_id": "movie:2"}, {$set: {sous_titre: "Le passager"}} )
db.test.update ({"_id": "movie:2"}, {$unset: {country: 1}} )
\end{smallverbatim}

\vfill

Et on supprime tout\,:

\begin{smallverbatim}
db.test.remove()
\end{smallverbatim}

\end{frame}


\section{Manipulation de données avec \mongodb/}



\begin{frame}[fragile]
\frametitle{Les requêtes}

Recherche par clé (très rapide). 

\begin{minted}[baselinestretch=.9,
               fontsize=\small,
               framesep=2mm]{javascript}
db.artists.find({_id: "artist:99"})
\end{minted}

Remarquez l'utilisation de \texttt{findOne}. Avec un chemin.

\begin{minted}[baselinestretch=.9,
               fontsize=\small,
               framesep=2mm]{javascript}
db.movies.find({"actors._id": "artist:99"})
\end{minted}

Avec pagination
\begin{minted}[baselinestretch=.9,
               fontsize=\small,
               framesep=2mm]{javascript}
db.movies.find({"actors._id": "artist:99"}).skip(0).limit(2)
\end{minted}

Avec une \alert{projection}

\begin{minted}[baselinestretch=.9,
               fontsize=\small,
               framesep=2mm]{javascript}
db.movies.find({"actors._id": "artist:99"}, {"title": 1})
\end{minted}

Avec une expression régulière (pas de guillemets!)\,:

\begin{minted}[baselinestretch=.9,
               fontsize=\small,
               framesep=2mm]{javascript}
db.movies.find({"title": /^Re/}, {"title": 1})
\end{minted}

Par intervalle:

\begin{minted}[baselinestretch=.9,
               fontsize=\small,
               framesep=2mm]{javascript}
db.movies.find({"year": {$gte: 2000, $lte: 2005}, {"title": 1})
\end{minted}

\end{frame}



\begin{frame}[fragile]
\frametitle{Les requêtes, suite}

Quelques clauses ensemblistes 

\begin{minted}[baselinestretch=.9,
               fontsize=\small,
               framesep=2mm]{javascript}
db.artists.find({_id: {$in: ["artist:34", "artist:98", "artist:1"]}, 
      {"last_name": 1})
\end{minted}

Autres possibilités\,: \texttt{\$nin}, \texttt{\$all}, \texttt{\$exist}

\begin{minted}[baselinestretch=.9,
               fontsize=\small,
               framesep=2mm]{javascript}
db.movies.find({'summary': {$exists: true}}, {"title": 1})
\end{minted}

\end{frame}



\begin{frame}[fragile]
\frametitle{Et les jointures}

Il faut les effectuer du \alert{côté client}. 

\vfill
Exemple (en supposant un schéma uniquement avec références): tous les films
réalisés par Clint Eastwood.

\vfill

On charge dans une variable \alert{locale} l'artiste Clint Eastwood.

\begin{minted}[baselinestretch=.9,
               fontsize=\small,
               framesep=2mm]{javascript}
eastwood = db.artists.findOne({"first_name": "Clint", "last_name": "Eastwood"})
\end{minted}

Notez le \texttt{findOne}\,!

\vfill

Et on recherche ses films.

\begin{minted}[baselinestretch=.9,
               fontsize=\small,
               framesep=2mm]{javascript}
db.movies.find({"director._id": eastwood['_id']}, {"title": 1})
\end{minted}
\vfill

Revient à effectuer une jointure par boucle imbriquées avec l'application client\,: 
très peu efficace à grande échelle.


\end{frame}

\begin{frame}[fragile]
\frametitle{Comparaison des techniques de jointure}


\figSlideWithSize{jointure-serveur-client}{10cm}

\end{frame}


\begin{frame}[fragile]
\frametitle{Une jointure complète}

On veut afficher les films et leur metteur en scène. Ce que l'on obtient en SQL avec:

\begin{minted}[baselinestretch=.9,
               fontsize=\small,
               framesep=2mm]{sql}
select m.titre, a.*  from Movie m, Artist a
where m.id_director = a.id
\end{minted}

\vfill
Voici le code MongoDB.

\begin{minted}[baselinestretch=.9,
               fontsize=\small,
               framesep=2mm]{javascript}
var lesFilms = db.movies.find()
while (lesFilms.hasNext()) {
  var film = lesFilms.next();
  var mes = db.artists.findOne({"_id": film.director._id});
  printjson(film.title);
  printjson(mes);
}
\end{minted}

\vfill
La requête imbriquée est exécutée autant de fois qu’il y a de films.

\end{frame}
